\section{Cubic response: adjoints, boundaries and Green kernels}\label{app:cubic-technical}
The conventions for $w$, $W$ and $J$ are those of \S\ref{sec:stationary}.
% Included at the end of Section 6, after the complete observable construction.
\subsection{Adjoint reduction of the second-order-field propagation products}\label{sec:prop-adjoint}
The part of the second-order ray driven by $H_2$ can be eliminated exactly. This subsection isolates \emph{all terms linear in the second-order fields}
$H_2$ and $\nu_2$, excluding the top-order fields $H_3,\nu_3$ and the remaining products made
only of first-order fields. It is a decomposition of the already complete map, not an omission
or a new initial-state assumption. The adjoint below holds at general momenta and finite times.

Write $h=H_1$, $j=H_2$, $m=\nu_2$, $P=1-nn^T$, $e=\bar t_\star$, $o=\bar t_o$,
$\bar X=n(o-t)$. In this subsection $h',j'$ denote partial time derivatives before pullback;
derivatives of ray and adjoint variables are total derivatives along $\bar X$.
Spatial derivatives act only on the fields. Let $X=X_1$, $v=v_1$ be the known first-order
orthonormal ray, and let $X_j,v_j$ denote only the part of $X_2,v_2$ linear in $j$.
Since $\zeta^{(2)}=0$, their equations and observer data are
\begin{equation}\label{eq:prop-linear-ray}
 X_j'=v_j+\tfrac12jn,\qquad
 v_j'=G_j=\tfrac12Pj'n+\tfrac12[\nabla j_{nn}-(n\cdot\partial)(jn)],\qquad
 X_j(o)=v_j(o)=0.
\end{equation}
The lapse $m$ does not deflect conformal null rays; it enters this sector only through the
physical emission surface and observer clock. Put
\begin{equation}\label{eq:prop-bilinear}
 \mathfrak B[h,j]=\int_e^o\left[
 \tfrac12X\cdot\nabla j'_{nn}-j'(n,v)
 +\tfrac12X_j\cdot\nabla h'_{nn}-h'(n,v_j)\right]\dd t.
\end{equation}
The cubic bulk contribution linear in $j$ is $3\mathfrak B[h,j]$ in the full-field Taylor
convention. Kernel polarisation instead gives the sum over its three single--pair partitions,
with no extra factor three.

\begin{proposition}[Propagation adjoint]\label{prop:propagation-adjoint}
Define first-order-only functions
\begin{gather*}
 r=\tfrac12\nabla h'_{nn},\qquad s=-Ph'n,\qquad \boldsymbol Y=v+\tfrac12P\boldsymbol K,\\
 \boldsymbol R(t)=\int_e^t r(u)\dd u,\qquad
 \boldsymbol K(t)=\int_e^t[s(u)-\boldsymbol R(u)]\dd u.
\end{gather*}
Then the unknown second-order ray can be eliminated exactly:
\begin{equation}\label{eq:prop-adjoint-first}
 \mathfrak B[h,j]=\int_e^o\left[
 \tfrac12X\cdot\nabla j'_{nn}-j'(n,v)
 -\tfrac12\boldsymbol R\cdot jn-\boldsymbol K\cdot G_j\right]\dd t.
\end{equation}
A further integration by parts removes all partial time derivatives of $j$ from the bulk:
\begin{equation}\label{eq:prop-adjoint-bulk}
\boxed{\begin{aligned}
 \mathfrak B[h,j]={}&
 [\tfrac12X\cdot\nabla j_{nn}-\boldsymbol Y\cdot jn]_e^o\\
 &+\int_e^o\left\{
 \tfrac12X\cdot\nabla(n\cdot\partial)j_{nn}
 -\tfrac12(X'+P\boldsymbol K)\cdot\nabla j_{nn}\right.\\
 &\hspace{19mm}\left.-v\cdot(n\cdot\partial)(jn)
 +(\boldsymbol Y'-\tfrac12\boldsymbol R)\cdot jn\right\}\dd t.
\end{aligned}}
\end{equation}
Only first-order ray and adjoint functions and the second-order metric remain. In particular,
$\boldsymbol R$ must not be projected transversely in its coupling to $jn$: the coordinate
displacement has a longitudinal part even though $v_j$ is transverse.
\end{proposition}
\begin{proof}
Differentiate $\boldsymbol R\cdot X_j+\boldsymbol K\cdot v_j$ and use
\eqref{eq:prop-linear-ray}. Its derivative is
$r\cdot X_j+s\cdot v_j+\boldsymbol R\cdot jn/2+\boldsymbol K\cdot G_j$.
Its endpoint values vanish, because the adjoints vanish at $e$ and the ray responses at $o$.
This proves \eqref{eq:prop-adjoint-first}. In the remaining $j'$ terms use
$\partial_tj=\dd j/\dd t+(n\cdot\partial)j$ on the background ray. The longitudinal
$\boldsymbol K$ pieces cancel between the two spatial-derivative terms, leaving
\eqref{eq:prop-adjoint-bulk}. The resulting metric boundary is \emph{not} zero:
at $o$, $X=v=0$ but $\boldsymbol Y=P\boldsymbol K/2$ need not vanish.
\end{proof}

\paragraph{Emission, physical clock and normalisation.}
Define $F_h=h'_{nn}/2$, $F_j=j'_{nn}/2$,
$U_h=\int_e^oF_h\dd t$, $U_j=\int_e^oF_j\dd t$ and
\[
 \mathcal C[m]=\frac{-\bar t_I^2m(\bar t_I,0)-\int_{\bar t_I}^otm(t,0)\dd t}{o}.
\]
This is the part of $D_2$ linear in $m$, not an additional clock correction.
The complete composed-exponent contribution of this sector is
\begin{equation}\label{eq:prop-complete-sector}
 \widehat I_3\big|_{2}=3\mathfrak D_2,\qquad
 \mathfrak D_2=\mathfrak B[h,j]
 +[eF_hm-F_jd_1]_e+D_1\partial_oU_j+\mathcal C[m]\partial_oU_h.
\end{equation}
Here $d_1=-e\nu_1(e,\bar X_e)$, and $D_1$ is the original physical-clock coefficient.
The $\partial_o$ derivatives include all background-ray and upper-limit dependence, with
$e$ fixed. After centring, the corresponding full normalised-map contribution is
\begin{equation}\label{eq:prop-normalized-sector}
 \Theta^{(3)}\big|_2=3\mathsf P_0\left[
 -\mathfrak D_2+(U_h-\sky{U_h})(U_j-\sky{U_j})\right].
\end{equation}
Thus the second-order-map normalisation product cannot be discarded. The dipole is still
retained in that product. All terms made only of first-order fields, including the mixed
clock--emission product in \eqref{eq:cubic-composed}, remain outside this isolated sector and
must be added as before. The Fermi--Walker correction is also first-order-field-only at this
map order; it does not add a $j$-dependent term here.

\paragraph{An evaluated shell emission contraction.}
The emission part of this sector already fits the nine-moment construction. Let
$\mathcal C_S(t)=\langle q_SF_S+q_TF_T\rangle_\mu$ (so $\Upsilon=\mathsf a\mathcal C_S$),
$Q_k=\hat k\hat k-\delta/3$. After polarising the full cubic expression and averaging over
$\hat q$, the emission-exponent kernel is exactly
\begin{equation}\label{eq:prop-emission-moments}
 \overline{\mathcal I_{3,e}\big|_2}(k;n)
 =\frac e2\left[
 \mathsf b'(n_h+3N_2-N_0)+3\mathsf a\mathcal C_S'
 \right]_e Q_{k,nn}\,e^{ik\cdot n(o-e)}.
\end{equation}
Indeed the two mixed partitions share one emission time, so their phases combine to the
external phase. The azimuthal averages are
$\overline{Q_q}=P_2(\mu)Q_k$ and
$\overline{H_2(k,q)}=\tfrac32(q_SF_S+q_TF_T)Q_k$.
The homogeneous tensor averages to zero and the homogeneous lapse remains.
For $\ell\ge2$, its temperature-transfer contribution is
\[
 \mathcal A_{\ell,e}\big|_2
 =\frac e2 w_\ell(k(o-e))
 [\mathsf b'(n_h+3N_2-N_0)+3\mathsf a\mathcal C_S']_e.
\]
The sign follows from $\Theta=-I$ at linear order in an exponent kernel and the
$Q_{k,nn}$ Funk--Hecke projection. This contribution can be added to the explicit emission
boundary of \eqref{eq:adjoint-nine-moment}; no identity setting their sum to zero is asserted.

\paragraph{What this does and does not settle.}
The bulk reduction and complete sector \eqref{eq:prop-complete-sector} are checked exactly
in \path{solution/check_propagation_adjoint.py}: matrix-exponential driver, arbitrary
plane-wave jets, ray duality, boundary and clock coefficients, centring algebra and shell
emission projections, with omission controls. It does not evaluate the full propagated
$13$ response. In a mixed-time product, for example,
\[
 (k+q)\cdot n(o-t)-q\cdot n(o-u)
 =k\cdot n(o-t)+q\cdot n(u-t).
\]
The second phase is present unless $u=t$. Consequently the finite nine polynomial moments
proved for the top-order fields and shared-time emission terms cannot simply be assumed to
close the remaining propagation integrals. No cancellation of the full cubic response or
late-time limit is proved. The unevaluated remainder is the combination of
\eqref{eq:adjoint-nine-moment}, \eqref{eq:prop-adjoint-bulk}, the clock terms and the remaining
first-order-only products, with the displayed physical boundaries retained.

\subsection{Field-reduced stationary functional}
\paragraph{Field-reduced form.} Eliminating the rays, endpoints, clock, frame and normalisation
into an exact reduced observable, and keeping as field rows the lapse equation (N), the five
orthonormal trace-free rows $2e^{H/2}(T)e^{-H/2}$, and the initial row
$\lim_{t\to0}\nu/t^2-c[g]$, gives
\begin{equation}\label{eq:stationary-field}
\begin{split}
 J_{abc}={}&J_{111,abc}
 +\int_0^o\dd t\Big[\Lambda_{N,0}C_{N,abc}+\Lambda_{H,0}:C_{H,abc}
 +\sum_{a|bc}\big(\Lambda_{N,a}C_{N,bc}+\Lambda_{H,a}:C_{H,bc}\big)\Big]\\
 &-\Lambda_{c,0}c_{abc}-\sum_{a|bc}\Lambda_{c,a}c_{bc}.
\end{split}
\end{equation}
Here $J_{111}$ is the complete measurement on the first-field-only family
$(\zeta,\nu,H)=\eps(g,a,h)$, $C_{N,H}$ are the direct residual derivatives along that family, and
$c_{bc},c_{abc}$ are the report's initial kernels (\S\ref{sec:fields}). There are three
partitions, each with coefficient one; the background adjoint uses the total momentum and the
first adjoint the pair momentum $k_b+k_c$, which may vanish. The Hamiltonian and momentum
constraints are \emph{not} added as further multiplier rows: they are consequences
(Thm.~\ref{thm:consistency}), and adding them would make $A$ redundant
(\S\ref{sec:constraint-rows} instead \emph{replaces} rows).

\paragraph{Background transpose and the regular-origin covector.} At momentum $K$ the inverse
of $L_K\nu=f_N$, $L_HH+4Q_K\nu=f_H$ is $\nu=b_cv_K+\int_0^tG_N^+f_N$ and
$H=\int_0^t\mathfrak G_K[f_H-4Q_K\nu]$, with $v_K=\frac9{K^2}u_1(Kt/\sqrt3)$ and $G_N^+$ the
retarded lapse kernel (opposite sign to the coefficient of $c_N$ in \S\ref{sec:fields}). Its
transpose has an initial-data covector that does not vanish:
\begin{equation}\label{eq:origin-covector}
 \Lambda_c=9\lim_{s\to0}s\Lambda_N(s),
\end{equation}
from $\lim_{s\to0}9sG_N^+(t,s)=v_K(t)$. For the linear quadrupole objective,
$J_y[\delta H]=-\frac12\int_e^oC_K:\delta H'\dd t$ with
$C_K(t)=\int\dd\Omega\,W\,\TF(nn)e^{iK\cdot n(o-t)}$, and
\[
 \Lambda_{H,0}(s)=\tfrac12\int_{\max(s,e)}^o\partial_t\mathfrak G_K(t,s)[C_K(t)]\dd t,\quad
 \Lambda_{N,0}(s)=-4\int_s^oG_N^+(t,s)Q_K:\Lambda_{H,0}(t)\dd t .
\]
The derivative acts on the kernel before any exchange of integrals, so the observer and emission
boundaries are kept. The nonlinear initial matching contributes
$-\Lambda_{c,0}c_{abc}-\sum\Lambda_{c,a}c_{bc}$ in \eqref{eq:stationary-field}; dropping these rows
loses part of the physical answer (see the gradient check below).

\paragraph{Exact normalisation adjoint.} With $F_Z=Z-e^{-I}$, $F_M=M-\sky{Z}$,
$F_\Theta=M(1+\Theta)-Z$ and $\int W\,\dd\Omega=0$, the multiplier of the exponent is
\begin{equation}\label{eq:normalisation-adjoint}
 \lambda_I=(1+\Theta)\Big[W-\frac{J}{4\pi}\Big],\qquad
 \lambda_{I,a}=W\Theta_a-\frac1{4\pi}\int W\Theta_a\,\dd\Omega,
\end{equation}
so $\int\lambda_I\,\dd\Omega=0$ exactly, and the random monopole has already cancelled from the
first adjoint source. The emission and clock multipliers have first variations
$\lambda_{{\rm em},a}=-e\,WF_a(e)$ and $\lambda_{{\rm clock},a}=-\partial_oJ^c_a/(a_1o)$; together
they reproduce the second-order clock functional
$\mathcal C[m]=-[\iota^2m(\iota,0)+\int_\iota^otm(t,0)\dd t]/o$.

\paragraph{The constant-direction source identities.}
For the dilation channel of \eqref{eq:dilation-channel}, the direct residuals satisfy
$C_{p00}=-2C_{p0}$, $c_{p00}=-4p^2/9=-2c_{p0}$ and $C_{00}=0$.
The reduced constant-direction adjoints are $-p\cdot\partial_p$ of the background ones.
These identities fix an exact zero-momentum channel, not a finite-radius shell average.

\paragraph{A physical check.} Applied to the leading-gradient map
$\Theta_{\rm lead}=-\tau^2e^{-2\eps g}[\eps\mathsf A-\eps^2\mathsf B]$ of \S\ref{sec:sh}, with the
nonlinear seed as a residual, the identity gives $\Theta^{(3)}_{\rm lead}=-12\tau^2(g^2\mathsf A+g\mathsf B)$,
in agreement with direct differentiation. Omitting the seed residual loses part of the
$g^2\mathsf A$ term.

\subsection{The complete first field-adjoint source}\label{sec:first-source}
The first field adjoint solves $A^*\Lambda_a=\mathcal S_a$, with
$\mathcal S_a[b]=-D_aJ_{{\rm red},y}[b]-\langle\Lambda_0,D_aE_y[b]\rangle$, on an independent
\emph{state} probe $b=(0,m(t),j(t))e^{iq\cdot X}$ (no primordial part). For the first input
$(\zeta,\nu,H)=(z,a,h)e^{ip\cdot X}$, $G=2z\mathbf 1+h$, $P=p+q$, $d=p\cdot q$, the off-shell
mixed residuals are
\begin{equation}\label{eq:mixed-residuals}
\begin{split}
 C_N&=\big[-q^ThP-z(2q^2-d)+4da+12a/t^2\big]m-a\,p^TjP+\tfrac12h':j',\\
 C_H&=2a'j'+2m'h'+2\TF\big[\mathcal R_{p,q}(G,j)-2\mathcal U-2am(pq^T+qp^T)\big],
\end{split}
\end{equation}
with $\mathcal U=\frac12\{m[P(Gq)^T+(Gq)P^T-dG]+a[P(jp)^T+(jp)P^T-dj]\}$ and
$\mathcal R_{p,q}$ the finite bilinear orthonormal Ricci coefficient (an explicit polynomial in
the momenta; no denominators). There are no second time derivatives of the probe. Two exact
controls: a homogeneous lapse probe ($q=0$, $j=0$) gives $C_N=12am/t^2$, $C_H=2m'h'$; a constant primordial shift
gives $C_N=-2zq^2m$, $C_H=-2z\mathfrak B_qj-8zQ_qm$, the finite-wavelength dilation $e^{-2z}$.

\paragraph{Measurement Hessian without a probe ray.} Let $\phi_q=e^{iq\cdot n(o-t)}$,
$N_n=nn^T-\frac13$, $F_a=\frac12\phi_ph'_{nn}$, $U_a=\int_e^oF_a$, and let $X_a,v_a$ be the known
fixed-coordinate-observer first ray, with propagation adjoints
$R_a=\int_e^t\frac12\nabla(\phi_ph'_{nn})$, $K_a=\int_e^t(s_a-R_a)$, $s_a=-P_n\phi_ph'n$
(as in Prop.~\ref{prop:propagation-adjoint}). With $\lambda_{I,a}$ from
\eqref{eq:normalisation-adjoint}, the first clock shift $D_a=\mathcal C[a]$ and the sky-integrated
$A_a=\int W\,\partial_oU_a\,\dd\Omega$,
\begin{equation}\label{eq:measurement-source}
\begin{split}
 S_{{\rm meas},a}[m,j]={}&\int_e^o\!\dd t\!\int\!\dd\Omega\,\phi_q\Big\{WE_0:j
 +\big[WE_1+\tfrac12\lambda_{I,a}N_n+\tfrac i2WD_a(q\cdot n)N_n\big]:j'\Big\}\\
 &+\int\dd\Omega\,\phi_q(e)\Big[e\,WF_a(e)m(e)+\tfrac e2Wa_eN_n:j'(e)\Big]
 +\frac{D_a}2\int\dd\Omega\,WN_n:j'(o)\\
 &-\frac{A_a}{o}\Big[\iota^2m(\iota)+\int_\iota^otm(t)\dd t\Big],
\end{split}
\end{equation}
where $E_1=\frac i2(q\cdot X_a)N_n-\TF\operatorname{sym}(v_an^T)-\frac12\TF\operatorname{sym}((P_nK_a)n^T)$
and $E_0=-\frac12\TF\operatorname{sym}(R_an^T)-\frac i2(q\cdot K_a)N_n
+\frac i2(q\cdot n)\TF\operatorname{sym}(K_an^T)$. The probe ray has been removed by the same
duality as in \S\ref{sec:prop-adjoint}; its total derivative vanishes only by
$R_a(e)=K_a(e)=0$ and $X_b(o)=v_b(o)=0$. The clock line lies on the observer worldline and is not
integrated over the sky again. The complete source is
$\mathcal S_a=S_{{\rm meas},a}-\int_0^o(V_\nu m+D_\nu m'+V_H:j+D_H:j')\dd t$, where
$V_\nu,D_\nu,V_H,D_H$ are the finite weak transposes of \eqref{eq:mixed-residuals} against
$\Lambda_{N,0},\Lambda_{H,0}$ at momentum $P$.

\paragraph{Frame.} With $R=e^{g(0)}e^{H(t,0)/2}E_{\rm FW}$, $R'=\Omega R$ and
$\Omega_{ab}=\frac18([h_a,h_b']+[h_b,h_a'])$, the first rotation vanishes, so the frame does not
source the first \emph{field} adjoint. Its first adjoint is
$\Lambda_{R,a}=-\mathcal K_a$, $\mathcal K_a=\int\dd\Omega\,\operatorname{skew}[W(\nabla_{S^2}\Theta_a)n^T]$,
and the cubic frame pairing $\sum_{\rm cyc}\int_\iota^o\mathcal K_a:\Omega_{bc}\dd t$ is retained.
It vanishes on the monochromatic shell, where $[h_a,h_b']=0$, but not for general triples.

\paragraph{Regular reconstruction and its boundaries.} On the causal Green probes
$\Phi_H=(0,\mathfrak G_q[B])$, $\Phi_N=(G_N^+,-4\int\mathfrak G_qQ_qG_N^+)$ and the origin probe
$\Phi_c$, the first adjoint is $\Lambda_{H,a}(s):B=\mathcal S_a[\Phi_H]$,
$\Lambda_{N,a}=\mathcal S_a[\Phi_N]$, $\Lambda_{c,a}=\mathcal S_a[\Phi_c]$. The point evaluations in
\eqref{eq:measurement-source} force
\begin{equation}\label{eq:first-adjoint-boundaries}
 \Lambda_{H,a}(o^-)=\frac{D_a}2\int WN_n\,\dd\Omega=\frac{D_a}{2}\mathsf C,\qquad
 \Lambda_{H,a}(e^-)-\Lambda_{H,a}(e^+)=\frac e2\int\dd\Omega\,Wa_e\phi_q(e)N_n,
\end{equation}
and $\Lambda_{c,a}=9\lim_{s\to0}s\Lambda_{N,a}(s)$. Imposing zero terminal data or a continuous
spatial adjoint would lose the physical clock and emission terms.

\subsection{The direct first-field measurement without a second ray}
For the first-field-only family write $v_2=u_2+n|v|^2$ with $n\cdot u_2=0$ (the radial part is
fixed by normalisation). Then $X_2'=u_2+Q_X$ and $u_2'=Q_v$, where
$Q_X=n|v|^2-Gv-\frac14G^2n+(X\cdot\nabla G)n$ and $Q_v$ is the explicit transverse forcing built
from first fields and first rays (working note \path{stationary_cubic_assembly.md}, eq.~(6)).
The same adjoints $R,K$ give the exact direct cubic line of sight
\begin{equation}\label{eq:direct-U3}
\boxed{U_3=\int_e^o\big[L_{111}-3R\cdot Q_X-3K\cdot Q_v\big]\dd t
 -3K(o)\cdot R_{{\rm FW},2}\,n,\qquad
 R_{{\rm FW},2}=\tfrac14\int_\iota^o[h(t,0),h'(t,0)]\dd t,}
\end{equation}
with the local remainder
$L_{111}=\frac18[h,[h,h']]_{nn}+\frac32X^iX^j\partial_i\partial_jh'_{nn}-6v^T(X\cdot\nabla h')n
+3v^Th'v-3|v|^2h'_{nn}$. No second ray remains; the frame enters through the observer value
$u_2(o)=-R_{{\rm FW},2}n$ and must not be added a second time through $\Lambda_R$. With the physical
emission shifts $d_1=-ea_e$, $d_2=e[a_e^2-2X_e\cdot\nabla a_e-2d_1(D_ta)_e]$, the clock shift
$D_1=\mathcal C[a]$ and
\begin{equation}\label{eq:second-clock}
 D_2=\frac{2\iota^3a(\iota,0)a'(\iota,0)-\int_\iota^ota(t,0)^2\dd t-2oa(o,0)D_1-D_1^2}{o},
\end{equation}
the exponent derivatives $I_{1,2,3}$ follow by composition (including
$-6D_1[(\partial_oF_{1,e})d_1+F_{1,e}\partial_od_1]$), and the exact normalised map is
$\Theta^{(3)}_{111}=P_0[-I_3+3ub-u^3+3\sky{u^2}u]$ with $u=I_1-\sky{I_1}$, $b=I_2-\sky{I_2}$.
Along the same family $D^3\operatorname{tr}S^2=0$, and the temporal part of the direct cubic
spatial residual is $\frac14[h,[h,h'']]+[h',[h,h']]+\frac1{2t}[h,[h,h']]$. \emph{On the
monochromatic shell every temporal commutator here and in \eqref{eq:direct-U3} vanishes
pointwise}, so the direct three-first-field Einstein source has no time derivatives of first
fields; spatial nonlinearities and transport remain.

\paragraph{The common fourth time.} The only four-time term,
$-3\int_e^o(n\cdot R)|v|^2\dd t$, has its common time integrated exactly (its integrand depends on
it only through the integration domain):
\[
 U_{3,R;abc}=-\sum_{S_3}\int_e^o\!\dd w\!\int_e^w\!\dd v\!\int_e^v\!\dd u\,(v-u)\,
 (n\cdot r_a(u))\,[V_b(v)\cdot V_c(w)]\,e^{i[p_a(o-u)+p_b(o-v)+p_c(o-w)]\cdot n}.
\]
Every term of the direct cubic measurement then has at most three independent times.

\paragraph{Independent measurement check.} In a spatially homogeneous anisotropic off-shell
metric, canonical photon momentum is conserved and
$e^{-I(n)}=[n^TR^Te^{H_o/2}e^{-H_e}e^{H_o/2}Rn]^{-1/2}$ at the physical roots. Expanding this
exactly reproduces \eqref{eq:direct-U3}--\eqref{eq:second-clock} and, with prescribed nonzero
$H_2,H_3,\nu_2$, the assembly $J_{\rm meas}^{(3)}=J_{111}-3S_{{\rm meas},1}[\nu_2,H_2]+J_{\rm lin}[H_3]$
using the implemented first source. This tests the measurement algebra, not a radiation solution.

\subsection{Exact angular closure and evaluated Green blocks}\label{sec:angular-green}
\paragraph{Entire phase moments.} Let
$F_r(z)={}_0F_1(;r+\frac32;-\frac z4)/(2r+1)!!$, so $F_r(x^2)=j_r(x)/x^r$ but with no
division. For a monomial of degree $d=\sum m_i$,
\begin{equation}\label{eq:phase-moments}
 \sky{n_1^{m_1}n_2^{m_2}n_3^{m_3}e^{iK\cdot n}}
 =i^d\sum_{0\le r_i\le m_i/2}(-1)^{r}F_{d-r}(K\cdot K)
 \prod_{i=1}^3\frac{m_i!\,K_i^{m_i-2r_i}}{2^{r_i}r_i!\,(m_i-2r_i)!},\qquad r=\textstyle\sum r_i,
\end{equation}
valid for all phase vectors, including $K=0$. Substituting the Volterra forms of the first ray
and adjoints, the \emph{entire measurement part} of the quadrupole first source becomes finite
combinations of $F_0,\dots,F_7$ at the mixed-time phase $K_{tu}=(p+q)(o-t)+p(t-u)$, with no
sky integral and no independently solved ray. The normalisation contribution contains, besides
the combined-phase term, the two separated products
$\frac14\mathcal U[h'_{nn};K_p]\mathcal A_C[N_n;K_q]+\frac14\mathcal A_C[h'_{nn};K_p]\mathcal U[N_n;K_q]$
($\mathcal U$ the plain and $\mathcal A_C$ the $w$-weighted phase average); both are generically
nonzero. Mixed-time phases mean this is \emph{not} a closure on the nine equal-time moments of
\S\ref{sec:cubic-moments}.

\paragraph{The frame covector is a commutator.} Since rotations preserve area,
$\mathcal K_a=[S_a,\mathsf C]$ with $S_a=\frac{15}2\sky{\Theta_aN_n}$ the linear quadrupole tensor;
for a scalar mode $S_a=-\frac{15}2T_2(|p|)Q_{\hat p}$. The same tensor gives $J_a=\mathsf C:S_a$ and
$A_a=-\mathsf C:\partial_oS_a$.

\paragraph{The forced homogeneous tensor drops out after isotropic contraction.} The exact pair
residuals at $(\alpha,-\alpha)$ are $f_H(t,|\alpha|)Q_\alpha$ and $f_N(t,|\alpha|)$, so on the regular
branch $H_{bc}=g_HQ_\alpha$. It enters the cubic response linearly, through a functional independent
of $\hat\alpha$; since $\ens{Q_\alpha}_{\hat\alpha}=0$, it cancels from the \emph{entire} isotropically
contracted cubic response, including measurement and boundary terms. The homogeneous lapse and
its initial row $c_{(q,-q)}=5q^2/9$ do not cancel. The uncontracted map still contains the tensor.

\paragraph{Evaluated Green kernels.} Let $\mathsf T_j(a,d)=\sum_{m\ge0}(-a)^md^{2m+j}/(2m+j)!$ (entire;
$\mathsf T_0,\mathsf T_1$ are the entire cosine and sine-over-frequency). On STF matrices
$\mathfrak B_q^3=\frac{2z}3\mathfrak B_q^2+\frac{z^2}3\mathfrak B_q$ ($z=q\cdot q$), and the retarded
spatial inverse of $\partial_t^2+\frac2t\partial_t+\mathfrak B_q$ is the quadratic operator polynomial
\begin{equation}\label{eq:spatial-green-poly}
 \mathfrak G_q(t,s)=\frac st\Big[d\,\mathbf 1-\frac{\mathsf T_3(z,d)+3\mathsf T_3(-\frac z3,d)}4\mathfrak B_q
 +\frac{3\mathsf T_5(z,d)+\mathsf T_5(-\frac z3,d)}4\mathfrak B_q^2\Big],\qquad d=t-s,
\end{equation}
which is the closed form of the series in Prop.~\ref{prop:cartesian}. With $a=z/3$,
$G_N^+=\frac13[\mathsf T_1(a,d)+(d\mathsf T_2-\mathsf T_3)(a,d)/(ts)]$, and the nested
lapse-to-spatial convolution is evaluated:
\[
\int_s^t\frac ut\mathsf T_1(-a,t-u)G_N^+(u,s)\dd u=\frac{s\mathcal A+\mathcal B+\mathcal D/s}{3t},
\]
with $\mathcal A=\frac12[\mathsf T_3(a)+\mathsf T_3(-a)]$, $\mathcal B=\frac12[d\mathsf T_3(a)+\mathsf T_4(-a)-\mathsf T_4(a)]$,
$\mathcal D=\frac12[d\mathsf T_4(a)-2\mathsf T_5(a)+\mathsf T_5(-a)]$ (all at $d$). The complete proper-clock
pairing is evaluated too. For a physical first input, set $a=p^2/3$ in the clock formulas:
$D_p=\frac{p^2}{9o}W_c$, with
$W_c=P_c(o)-P_c(\iota)+\iota^2v_p(\iota)$ and
$P_c(T)=3[T\mathsf T_3(a,T)-2\mathsf T_4(a,T)]$; at zero momentum
$W_c=(o^4+3\iota^4)/4$.

For the physical first scalar mode, the direct primordial spatial forcing
$-2Q\mathsf T_3(-a,t)/t$ and the initial-row column $c_1H_c$ ($c_1=-a/3$) combine to
$H_1=-2Q[\mathsf T_2(a,t)-2\mathsf T_3(a,t)/t]$: all hyperbolic ($-a$) terms cancel between the source and
the initial matching. \S\ref{sec:constraint-rows} shows why, and extends this to every physical
pair.

\paragraph{Time arity.} Counting independent integration variables after distributing products
(a product of three one-time means counts as three): background adjoints need one time, first
adjoints at most two, the single--pair pairings and the direct cubic measurement at most three,
initial rows at most two. Hence the complete stationary cubic \emph{quadrupole} has at most three
independent time integrations per distributed term, for arbitrary triples. This bounds the
representation; it does not evaluate the surviving kernels.

\subsection{The constraint-row transpose}
\paragraph{The transpose, with its boundary terms.} Let a weak spatial functional be
$\mathcal W[j]=\int_0^o(\Omega_0:j+\Omega_1:j')\dd t+\Omega_e:j'(e)+\Omega_o:j'(o)$, plus lapse terms. On
compatible sources at $K\neq0$, its non-TT part equals
\begin{equation}\label{eq:constraint-transpose}
\begin{gathered}
 \int_0^oZ(s)\cdot\big[4K(\nu'-\nu/s)+2ic_M\big]\dd s
 +\sum_{x=e,o}\mathcal P_K^T[\Omega_x]\cdot\big[4K(\nu'-\nu/t)+2ic_M\big]_{t=x},\\
 Z(s)=\mathcal P_K^T\Big[\Omega_1(s)+\int_s^o\Omega_0\Big],
\end{gathered}
\end{equation}
with $\mathcal P_K^T[\Omega]=\frac2{K^2}\Omega K-\frac{K\cdot\Omega\cdot K}{2K^4}K$. The lapse weights then
pass through the acoustic transpose and its origin covector \eqref{eq:origin-covector}. No
hyperbolic function occurs. For the linear quadrupole, $\Omega_1=-\frac12C_K\mathbf 1_{(e,o)}$ and the
momentum-row weight is the \emph{local} $Z_0=-\frac12\mathcal P_K^T[C_K(t)]$. This has the
structure of the momentum form of Prop.~\ref{prop:adjoint}: an acoustic lapse Wronskian plus a
local momentum-source term.

\paragraph{Application to the single--pair term.} Each pair source $D_bD_cF$ is a physical
second-order source and so is compatible. Since
$\langle\lambda_a,D_bD_cF\rangle=-\mathcal S_a[y_{bc}]$, the hyperbolic scalar part of the first adjoint
at the probe momentum $k_b+k_c\neq0$ drops out exactly: \eqref{eq:constraint-transpose} applied
to $\mathcal S_a$ replaces it by acoustic kernels, one cumulative integral and the explicit point
terms at $e$ and $o$ that carry \eqref{eq:first-adjoint-boundaries}. The \emph{background}
adjoint inside $\mathcal S_a$ (total momentum) is paired with the mixed cubic residual, which is
not compatible by itself. Only the sum of the direct third residual and the three field-variation
parts is the full physical cubic source, to which the same replacement applies
(Prop.~\ref{prop:combined}). The separate background terms of \eqref{eq:stationary-field} keep
hyperbolic parts that cancel only in that sum. The individual terms could be made hyperbolic-free
only by rewriting the whole stationary functional in the momentum rows. That would require their
mixed variations, which are not derived here and, by Prop.~\ref{prop:combined}, are not needed.

\subsection{Four-gradient first-adjoint coefficients}
\paragraph{The cubic quadrupole through four gradients.} Scale momenta by $\kappa$. The leading
first spatial adjoint is evaluated in closed form, with observer trace
$p^2(o^4+3\iota^4)\mathsf C/(72o)$ and emission jump $-p^2e^3\mathsf C/18$; its integral has no
logarithms,
\[
 \overline\Lambda_{H,a}^{[2]}=\frac L3U+\tau^4V-\frac E{120}\rho_{p,q}
 +\frac{p^2(3o^4+5\iota^4-8e^4)}{360}\mathsf C,\qquad
 \Lambda^{[2]}_{c,a}=\frac{3o^4+5\iota^4-8e^4}{20}\mathsf C:Q_p+\frac E5(p+q)^T\mathsf Cq,
\]
with $\tau^2=(o^2-e^2)/6$, $E=o^4-e^4$, $L=\int_e^ot(o-t)^2\dd t$, and $U,V,\rho_{p,q}$ explicit
quadratic tensors in $p,q,\mathsf C$ (note \path{quadrupole_adjoint_gradient_expansion.md}). Both
initial-data pairings enter at four gradients, below the bulk lapse terms.
These coefficients use the averaged objective $J=\sky{w\Theta}$ and its
equivalent full-solid-angle weight $W=w/(4\pi)$.
